\documentclass[12pt, a4paper]{article}
\usepackage{../../../myconfig} % Gọi file cấu hình từ ngoài gốc vào

% ==========================================
% BẮT ĐẦU NỘI DUNG TÀI LIỆU
% ==========================================
\begin{document}

% Truyền 3 tham số: {Nhãn cam} {Chữ to chính giữa} {Chữ ở góc phải Header}
\titlebox{Chủ đề: Hàm số}{Lượng giác}{Hàm số: Lượng giác}

\drawdashlines{20}
\newpage
\drawdashlines{25}
\newpage
% ==========================================
% KHỐI CÂU HỎI TRẮC NGHIỆM LƯỢNG GIÁC
% ==========================================

% --- CÂU 1 ---
\questionLinesManual{13}{1}{Nghiệm của phương trình \( \cos x = \dfrac{1}{2} \) là:}{}

% --- CÂU 2 ---
\questionLinesManual{13}{2}{Nghiệm của phương trình \( \cos x = \cos \dfrac{\pi}{12} \) là:}{}

% --- CÂU 3 ---
\questionLinesManual{13}{3}{Tập nghiệm của phương trình \( \sin x = 1 \) là:}{}

% --- CÂU 4 ---
\questionLinesManual{13}{4}{Họ nghiệm của phương trình \( \sin\left(x + \dfrac{\pi}{7}\right) = \dfrac{\sqrt{3}}{2} \) là:}{}

% --- CÂU 5 ---
\questionLinesManual{13}{5}{Tìm tập xác định \( D \) của hàm số \( y = \tan 2x \):}{}

% --- CÂU 6 ---
\questionLinesManual{13}{6}{Tìm tập xác định \( D \) của hàm số \( y = \dfrac{3\sin x + 1}{1 - \cos 2x} \):}{}

\newpage
\drawdashlines{18}

% --- CÂU 7 ---
\questionLinesManual{13}{7}{Số nghiệm của phương trình \( 2\cos x - 1 = 0 \) trên đoạn \([0; 2\pi]\) là:}{}

% --- CÂU 8 ---
\questionLinesManual{13}{8}{Số nghiệm của phương trình \( \sin x = -\dfrac{\sqrt{3}}{2} \) trên đoạn \(\left[-\dfrac{\pi}{2}; 4\pi\right]\) là:}{}

% --- CÂU 9 ---
\questionLinesManual{13}{9}{Số nghiệm của phương trình \( \dfrac{\sin 3x}{\cos x + 1} = 0 \) thuộc đoạn \([2\pi; 4\pi]\) là:}{}


\end{document}